\documentclass[11pt]{article}
\usepackage{amsmath,amssymb,amsthm,mathtools}
\usepackage[margin=25mm]{geometry}
\usepackage{hyperref}
\newcommand{\C}{\mathbb C}
\newcommand{\R}{\mathbb R}
\newcommand{\conj}[1]{\overline{#1}}
\title{Two boundary polynomials for the original Gamma kernel\\and the exact correction after coefficient compression}
\author{}
\date{20 September 2026}
\begin{document}\maketitle

\section{Source, mass, coordinates and scope}

This is a finite calculation on the original Gamma receiver MR22--23
\cite{MR}%
% BEGIN VERIFIED PUBLIC PROOF LINK
~\href{https://github.com/KokunoYumeto/zeta-function-research-reader/blob/5992ad50c0f2a06921f1fcaded7b4e38097c0739/workbenches/splitzero-tandem/continuations/20260920-original-kernel-mixed-determinants/ORIGINAL_KERNEL_MIXED_DETERMINANT.tex\#L36}{[proof MR1]}%
% END VERIFIED PUBLIC PROOF LINK
{}, not a replacement of the arithmetic source or its original
coefficient frame. The underlying Christoffel--Darboux identity is
classical; Akemann--Vernizzi's original author source explicitly gives
its real-line formula and its algebraic continuation in the footnote
at \texttt{CD} \cite{AV}. The exact Gamma recurrence and norms below
are the Meixner--Pollaczek formulas \cite{DLMF}, transported in the
original JM26 and KG6 coordinates \cite{JM,KG}%
% BEGIN VERIFIED PUBLIC PROOF LINK
~\href{https://github.com/KokunoYumeto/zeta-function-research-reader/blob/5992ad50c0f2a06921f1fcaded7b4e38097c0739/workbenches/splitzero-tandem/continuations/20260920-original-kernel-mixed-determinants/providers/JOINT_MINIMUM_RECEIVERS.tex\#L286}{[proof JM21]} \href{https://github.com/KokunoYumeto/zeta-function-research-reader/blob/5992ad50c0f2a06921f1fcaded7b4e38097c0739/workbenches/splitzero-tandem/continuations/20260920-original-kernel-mixed-determinants/providers/KERNEL_GAMMA_INDEPENDENT.tex\#L73}{[proof KG1]}%
% END VERIFIED PUBLIC PROOF LINK
{}. We derive every finite
identity used here. No large-degree coefficient, limit estimate,
arithmetic zero statement, or priority claim about the classical
identity is asserted.

The retained density, original physical coordinate and its complex
evaluation chart are
\begin{equation}\label{eq:source}
 d\sigma(y)=\frac{|\Gamma(1/4+iy/2)|^2}{2\pi}\,dy,
 \quad M_\sigma=\sqrt{2\pi},\quad c=k/2,
 \quad S=c+iy,
 \quad u=\frac{v-c}{i},\quad w=\frac{z-c}{i}.
 \tag{CD1}
\end{equation}
In formulas with two scalar polynomial variables, write \(x=u\) and
\(t=\overline w\); these variables are independent during
differentiation. An asterisk on a matrix always means conjugate
transpose. The source inner product is
\(\langle P,Q\rangle=\int\overline{P(c+iy)}Q(c+iy)\,d\sigma(y)\).

Here is the mass check, retaining the scale in \(y\). The original DLMF
weight and norm at \(\lambda=1/4,\phi=\pi/2\) are
\[
 |\Gamma(1/4+ix)|^2\,dx,
 \qquad
 h_j=\frac{2\pi\Gamma(j+1/2)}{\sqrt2\,j!}.
\]
Set \(y=2x\) and
\(p_j(y)=j!P_j^{(1/4)}(y/2;\pi/2)\).
The measure becomes \(|\Gamma(1/4+ix)|^2dx/\pi\). Hence the actual
squared norm is
\begin{equation}\label{eq:norm}
 \gamma_j=\frac{(j!)^2}{\pi}h_j
 =\sqrt2\,j!\Gamma(j+1/2)
 =\sqrt{2\pi}\,j!(1/2)_j.
 \tag{CD2}
\end{equation}
In particular \(\gamma_0=M_\sigma=\sqrt{2\pi}\). The exact density
change \(dy=2dx\) is already included in CD2.

The original generating function, after the same substitution, is
\[
 \sum_{j\geq0}\frac{p_j(y)}{j!}\zeta^j
 =(1+\zeta^2)^{-1/4}e^{y\arctan\zeta}.
\]
Its highest \(y\) coefficient proves monicity. Differentiating in
\(\zeta\) gives
\((1+\zeta^2)\partial_\zeta F=(y-\zeta/2)F\).
Coefficient comparison therefore proves the original recurrence
\begin{equation}\label{eq:recurrence}
 p_0=1,\quad p_1=y,\quad
 p_{j+1}(y)=yp_j(y)-a_jp_{j-1}(y),
 \quad a_j=j(j-\tfrac12)=\frac{\gamma_j}{\gamma_{j-1}}
 \quad(j\geq1).
 \tag{CD3}
\end{equation}
All coefficients of \(p_j\) are real. We put \(p_{-1}=0\) only in
recurrences whose coefficient is zero; no negative-index norm is used.

\section{The two-polynomial identity and its sign}

For \(d\geq0\), set
\[
 \mathcal K_d(x,t)=\sum_{j=0}^{d-1}\frac{p_j(x)p_j(t)}{\gamma_j},
 \qquad K_d(v,\overline z)=\mathcal K_d(u,\overline w).
\]
An empty sum is zero. For \(d\geq1\), define
\[
 F_d(x,t)=p_d(x)p_{d-1}(t)-p_{d-1}(x)p_d(t).
\]
Then the exact polynomial identity is
\begin{equation}\label{eq:cd}
 (x-t)\mathcal K_d(x,t)=\frac{F_d(x,t)}{\gamma_{d-1}}.
 \tag{CD4}
\end{equation}
To prove it without dividing by \(x-t\), put
\(T_j=p_{j+1}(x)p_j(t)-p_j(x)p_{j+1}(t)\).
The recurrence gives
\[
 (x-t)p_j(x)p_j(t)=T_j-a_jT_{j-1}.
\]
For \(j=0\) the second term is zero. Division by \(\gamma_j\) and
\(a_j/\gamma_j=1/\gamma_{j-1}\) make the finite sum telescope to
\(T_{d-1}/\gamma_{d-1}=F_d/\gamma_{d-1}\).
This proves \eqref{eq:cd} as an identity in \(\C[x,t]\), including
\(x=t\). Away from that diagonal it gives
\begin{equation}\label{eq:offdiagonal}
 \boxed{\mathcal K_d(x,t)=
 \frac{p_d(x)p_{d-1}(t)-p_{d-1}(x)p_d(t)}
      {\gamma_{d-1}(x-t)}.}
 \tag{CD5}
\end{equation}
At the diagonal, differentiation of the polynomial identity in \(x\)
and then setting \(x=t\) gives
\begin{equation}\label{eq:diagonal}
 \boxed{\mathcal K_d(t,t)=
 \frac{p_d'(t)p_{d-1}(t)-p_{d-1}'(t)p_d(t)}{\gamma_{d-1}}.}
 \tag{CD6}
\end{equation}
At \(d=1\) both formulas give \(1/M_\sigma\), fixing the denominator
sign. At \(d=0\), the kernel and all its derivatives are zero; neither
CD5 nor CD6 is used.

The collision in these formulas is \(u=\overline w\), not merely
\(u=w\). For complex off-line arguments, the Hermitian kernel entry
\(\mathcal K_d(u,\overline u)\) usually belongs to CD5, not CD6.
CD6 applies when the two polynomial arguments coincide, whether or not
that common argument is real. Positivity is asserted for Hermitian Gram
matrices, not for every analytically continued scalar entry.

\section{Every divided jet, both away from and at a collision}

For a polynomial \(f\), write \(f^{[a]}=f^{(a)}/a!\). Define
\[
 F_{d;a,b}(x,t)=p_d^{[a]}(x)p_{d-1}^{[b]}(t)
             -p_{d-1}^{[a]}(x)p_d^{[b]}(t),
 \qquad
 \mathcal K_{d;r,s}=\frac{\partial_x^r\partial_t^s\mathcal K_d}{r!s!}.
\]
For all nonnegative integers \(r,s\), at \(x\ne t\),
\begin{equation}\label{eq:jets-away}
 \boxed{\mathcal K_{d;r,s}(x,t)=\frac1{\gamma_{d-1}}
  \sum_{a=0}^r\sum_{b=0}^s
  \frac{(-1)^a\binom{a+b}{a}\,
        F_{d;r-a,s-b}(x,t)}{(x-t)^{a+b+1}}.}
 \tag{CD7}
\end{equation}
Indeed the divided \((a,b)\) derivative of \((x-t)^{-1}\) is
\((-1)^a\binom{a+b}{a}(x-t)^{-a-b-1}\). The divided-derivative product
rule is convolution with no further binomial factor. Applying it to
CD5 proves CD7.

At a collision \((x,t)=(\tau,\tau)\), the exact formula is
\begin{equation}\label{eq:jets-collision}
 \boxed{\mathcal K_{d;r,s}(\tau,\tau)
  =\frac1{\gamma_{d-1}}
      \sum_{j=0}^{s}F_{d;r+1+j,s-j}(\tau,\tau)
  =-\frac1{\gamma_{d-1}}
      \sum_{j=0}^{r}F_{d;r-j,s+1+j}(\tau,\tau).}
 \tag{CD8}
\end{equation}
For a complete proof, expand both sides of
\(F_d(x,t)=\gamma_{d-1}(x-t)\mathcal K_d(x,t)\) at
\(x=\tau+h,t=\tau+z\). Equality of the coefficient of \(h^az^b\)
gives
\[
 F_{d;a,b}(\tau,\tau)=\gamma_{d-1}
 \bigl(\mathcal K_{d;a-1,b}(\tau,\tau)
       -\mathcal K_{d;a,b-1}(\tau,\tau)\bigr),
\]
where any negative derivative index means zero. Starting with
\((a,b)=(r+1,s)\) and iterating downward in \(s\) proves the first
sum. Starting with \((r,s+1)\) and iterating downward in \(r\)
proves the second. The expansion is finite, so there is no interchange
of infinite limiting operations. CD7--8 include arbitrary repeated
roots and all mixed divided jets. Derivatives beyond the degree are
zero, exactly as in the original finite kernel.

\section{The original \(S\) coordinate and its reflection denominator}

The monic polynomials in the physical coordinate are
\[
 P_j(v)=i^jp_j((v-c)/i),\qquad \|P_j\|^2=\gamma_j.
\]
The paired phases cancel in each kernel summand, so this is the same
\(K_d\) already defined. Since
\[
 u-\overline w=-i(v+\overline z-2c),
\]
and
\[
 F_d(u,\overline w)=-i\bigl(
 P_d(v)\overline{P_{d-1}(z)}+
 P_{d-1}(v)\overline{P_d(z)}\bigr),
\]
CD5 becomes
\begin{equation}\label{eq:S-CD}
 \boxed{K_d(v,\overline z)=
 \frac{P_d(v)\overline{P_{d-1}(z)}+
       P_{d-1}(v)\overline{P_d(z)}}
      {\gamma_{d-1}(v+\overline z-2c)}.}
 \tag{CD9}
\end{equation}
This expression is used away from \(v+\overline z=2c\); the
polynomial extension supplies the remaining values. The numerator is a
sum and the denominator is a reflection sum. Replacing either by the
usual real-line difference would change the original coordinate.

The exact divided-jet transport is
\begin{equation}\label{eq:S-jets}
 \frac{\partial_v^r\partial_{\overline z}^sK_d(v,\overline z)}{r!s!}
   =i^{-r}i^s\mathcal K_{d;r,s}(u,\overline w).
 \tag{CD10}
\end{equation}
This follows from \(\partial_v=i^{-1}\partial_u\) and
\(\partial_{\overline z}=i\partial_{\overline w}\), so it also
fixes the phases at collisions. For a direct formula in physical
variables, write \(G_d(v,b)=P_d(v)\overline{P_{d-1}(\overline b)}+
P_{d-1}(v)\overline{P_d(\overline b)}\), a polynomial in independent
variables \(v,b\), and use divided derivatives \(G_{d;a,b}\), with
the second subscript here an integer. Away from \(v+b=2c\), the
counterpart of CD7 has denominator \((v+b-2c)^{a+b+1}\) and sign
\((-1)^{a+b}\) in each summand. At a point \(v_0+b_0=2c\), it is
\begin{equation}\label{eq:S-jets-collision}
 \frac{\partial_v^r\partial_b^sK_d(v,b)}{r!s!}
 \bigg|_{(v_0,b_0)}
  =\frac1{\gamma_{d-1}}
      \sum_{j=0}^s(-1)^jG_{d;r+1+j,s-j}(v_0,b_0).
 \tag{CD11}
\end{equation}
To prove this independently, expand
\(G_d=\gamma_{d-1}(v+b-2c)K_d\) at that point; the local linear
factor is now \(h+z\), so solving the same coefficient recurrence
alternates the signs. This agrees with CD10.

\section{Every entry at the original quartet and all four cutoffs}

The original simple-packet domain is
\[
 0<\delta<\tfrac12,\quad \gamma>2,\quad k\equiv1\pmod4,
 \quad k\geq9,\quad q=(k+1)^2,
\]
on the original five-orbit period locus of OCF \cite{OCF}%
% BEGIN VERIFIED PUBLIC PROOF LINK
~\href{https://github.com/KokunoYumeto/zeta-function-research-reader/blob/5992ad50c0f2a06921f1fcaded7b4e38097c0739/workbenches/splitzero-tandem/continuations/20260920-original-kernel-mixed-determinants/providers/ORIGINAL_CONDUCTOR_STRIP_FRAME.tex\#L27}{[proof OCF1]}%
% END VERIFIED PUBLIC PROOF LINK
{}. Here
\(\gamma\) is the retained zero height, whereas \(\gamma_j\) is a
squared polynomial norm. The ordered roots and their exact charts are
\begin{equation}\label{eq:roots}
 v_{ab}=c+(2a-k)\delta+i(2b-k)\gamma,
 \qquad u_{ab}=(2b-k)\gamma-i(2a-k)\delta,
 \quad 0\leq a,b\leq k.
 \tag{CD12}
\end{equation}
Distinctness follows by comparing the two real coordinates, since
\(\delta,\gamma>0\). However
\[
 u_{ab}=\overline{u_{a'b'}}
 \quad\Longleftrightarrow\quad b=b',\quad a'=k-a.
\]
Indeed equality of real parts gives \(b=b'\), and equality of imaginary
parts gives \(-(2a-k)=2a'-k\). Thus the full covariance has exactly
one collision entry per row, at its reflected column \((k-a,b)\).
There are \(q\) such ordered entries. Since \(k\) is odd, there are no
fixed points \(a=k-a\), and these are not its Hermitian diagonal
entries. Their CD6 values must be retained.

For \(d=N+1\), let \(C_d\) denote the matrix
\([\mathcal K_d(u_\alpha,\overline{u_\beta})]\) in exactly that
lexicographic root order. For every pair of entries it is computed from
the same two polynomials by
\begin{equation}\label{eq:entry-recipe}
 (C_d)_{\alpha\beta}=
 \begin{cases}
 \dfrac{p_d(u_\alpha)\overline{p_{d-1}(u_\beta)}-
 p_{d-1}(u_\alpha)\overline{p_d(u_\beta)}}
 {\gamma_{d-1}(u_\alpha-\overline{u_\beta})},
 &u_\alpha\ne\overline{u_\beta},\\[7pt]
 \dfrac{p_d'(u_\alpha)p_{d-1}(u_\alpha)-
 p_{d-1}'(u_\alpha)p_d(u_\alpha)}{\gamma_{d-1}},
 &u_\alpha=\overline{u_\beta}.
 \end{cases}
 \tag{CD13}
\end{equation}
This is an exact two-boundary-polynomial expression, including the
derivatives of those same polynomials at the resonant entries.
The four original cutoffs require exactly these pairs:
\[
\begin{array}{c|c|c|c}
N&d&\text{boundary pair}&\text{denominator norm}\\\hline
q-1&q&(p_q,p_{q-1})&\gamma_{q-1}\\
q&q+1&(p_{q+1},p_q)&\gamma_q\\
2q-1&2q&(p_{2q},p_{2q-1})&\gamma_{2q-1}\\
2q&2q+1&(p_{2q+1},p_{2q})&\gamma_{2q}
\end{array}
 \tag{CD14}
\]
Every number in the table follows from \(d=N+1\); no degree shift is
absorbed into a different kernel convention.

The original nonreal frame from OCF and MR6--8 is retained literally:
\[
 Z=[\,M_A\ \ S_kZ_k\,]\in\C^{q\times r},\qquad
 \pi Z=0,\quad r=k^2-6k+17,
 \quad m=q-r=8k-16.
\]
All its period and amplitude coefficients, pivot choices and column
order remain those of the original construction. The compressed matrix
is therefore given entry by entry by
\begin{equation}\label{eq:compression-recipe}
 A_d=Z^TC_d\overline Z,
 \qquad
 (A_d)_{ij}=\sum_{\alpha,\beta=1}^q
    Z_{\alpha i}(C_d)_{\alpha\beta}\overline{Z_{\beta j}},
 \tag{CD15}
\end{equation}
with CD13 substituted in every term. Thus two boundary polynomials
suffice for the full original compression at each cutoff, provided all
root entries, all nonreal frame coefficients and all collision values
are retained. This statement is not an assertion that the two already
compressed boundary vectors alone obey a closed rank-two displacement.
The exact additional term is calculated next.

\section{Full displacement and its confluent version}

Put \(Y=\operatorname{diag}(u_\alpha)\) and
\(f_j=(p_j(u_\alpha))_\alpha\). CD4 is precisely the matrix identity
\begin{equation}\label{eq:full-displacement}
 YC_d-C_dY^*=
 \frac{f_df_{d-1}^*-f_{d-1}f_d^*}{\gamma_{d-1}}
 =:B_d.
 \tag{CD16}
\end{equation}
The right side is skew-Hermitian and has rank at most two. For a
resonant entry the scalar multiplier on the left is zero, and the
corresponding numerator on the right is zero. Consequently CD16 alone
does not determine those entries of \(C_d\); CD6 or CD8 supplies them.
For example a matrix supported at a single resonant entry belongs to
the kernel of \(X\mapsto YX-XY^*\). This proves the need for the
collision data without assuming uniqueness for a singular matrix map.

For arbitrary finite divided-jet blocks, define instead
\[
 f_j=\bigl(p_j^{[a]}(u_\alpha)\bigr)_{(\alpha,a)},
 \qquad J=\bigoplus_\alpha(u_\alpha I+L_\alpha),
 \qquad (L_\alpha)_{a,a-1}=1.
\]
The other entries of each \(L_\alpha\) are zero; rows are ordered by
increasing derivative degree. Divided Leibniz differentiation gives
\((yp_j)^{[a]}(u_\alpha)=u_\alpha p_j^{[a]}(u_\alpha)
+p_j^{[a-1]}(u_\alpha)\). Hence
\(Jf_j=f_{j+1}+a_jf_{j-1}\), and the same telescoping proof gives
\begin{equation}\label{eq:jet-displacement}
 JC_d-C_dJ^*=
 \frac{f_df_{d-1}^*-f_{d-1}f_d^*}{\gamma_{d-1}},
 \quad C_d=\sum_{j<d}\frac{f_jf_j^*}{\gamma_j}.
 \tag{CD17}
\end{equation}
Its entries are exactly the divided derivatives in CD7--8. If the total
jet count is at most \(d\), the first that many monic polynomial columns
have full row rank: a polynomial of smaller degree in the kernel would
be divisible by the product of all the specified root powers, by its
finite Taylor expansions, and would therefore be zero. Thus \(C_d>0\)
in the original admitted-degree regime. The displacement identity itself
does not require this positivity or a matrix inverse.

\section{Exact compression and the complementary coupling term}

For the original simple-root frame put
\[
 V=\overline Z,\quad W=V^*V=Z^T\overline Z,\quad
 P=VW^{-1}V^*,\quad Q=I-P,\quad
 H=(V^*YV)W^{-1},\quad g_j=V^*f_j=Z^Tf_j.
\]
Full column rank of \(Z\), proved in OCF and MR7, makes \(W>0\).
Direct multiplication gives \(P^*=P\), \(P^2=P\),
\(PV=V\), and \(V^*Q=0\). Therefore \(P\) is the Euclidean
orthogonal projection to \(\operatorname{im}V\). This projection is
an algebraic device; the original source Gram and the original columns
of \(Z\) are not changed.

The placement of \(W^{-1}\) is essential: \(A_d=V^*C_dV\) is a
covariant Gram, so the reduced operator appearing on its left is
\(H=(V^*YV)W^{-1}\). The coefficient operator with the opposite
placement \(W^{-1}(V^*YV)\) cannot be inserted into the following
formula without also transporting the Gram.

Define the two actual complement maps
\[
 L=QY^*V,\qquad T_d^{\perp}=QC_dV,\qquad
 \mathcal E_d=L^*T_d^{\perp}-(T_d^{\perp})^*L.
\]
Then the exact compressed displacement is
\begin{equation}\label{eq:compressed-displacement}
 \boxed{HA_d-A_dH^*=
 \frac{g_dg_{d-1}^*-g_{d-1}g_d^*}{\gamma_{d-1}}
   -\mathcal E_d.}
 \tag{CD18}
\end{equation}
To prove it, use
\[
 HA_d=V^*YPC_dV,\qquad A_dH^*=V^*C_dPY^*V.
\]
Insert \(P=I-Q\), apply CD16 to the full terms, and retain the
remaining difference
\(V^*YQC_dV-V^*C_dQY^*V=L^*T_d^\perp-(T_d^\perp)^*L\).
This proves both the sign and every factor in CD18.

Equivalently, put \(\rho_j=V^*YQf_j\). Projecting the original
three-term recurrence gives its exact inhomogeneous receiver
\[
 Hg_j=g_{j+1}+a_jg_{j-1}-\rho_j.
\]
Since \(T_d^\perp=\sum_{j<d}Qf_jg_j^*/\gamma_j\), its correction is
\begin{equation}\label{eq:correction-sum}
 \mathcal E_d=\sum_{j<d}
    \frac{\rho_jg_j^*-g_j\rho_j^*}{\gamma_j}.
 \tag{CD19}
\end{equation}
Thus all lower-degree complement couplings occur, with their exact
signs, in addition to the two projected boundary vectors. Formula CD15
still computes this entire matrix from the two full boundary
polynomials and their collision data; CD19 describes what a closed
recurrence on compressed vectors alone would otherwise omit.

The original complementary frame is available without selecting a new
one. Put
\[
 U=\pi^T\in\C^{q\times m},\qquad M=U^*U=\overline\pi\pi^T.
\]
The coefficient identity \(\pi Z=0\) gives \(U^*V=0\); the ranks
sum to \(q\), so \(Q=UM^{-1}U^*\). Define
\[
 F=V^*YU,\qquad T_d=U^*C_dV.
\]
Substitution into the already proved correction gives
\begin{equation}\label{eq:actual-complement}
 \boxed{\mathcal E_d=FM^{-1}T_d-T_d^*M^{-1}F^*,
 \qquad
 \operatorname{rank}(HA_d-A_dH^*)\leq\min(r,2+2m).}
 \tag{CD20}
\end{equation}
Each correction summand has rank at most \(m\), so their difference
has rank at most \(2m\). Adding the rank-at-most-two boundary term
proves the bound. Equivalently the displacement's image lies in the
span of the \(2+2m\) columns
\([g_d,g_{d-1},F,T_d^*]\). The exact original dimensions are
\(r=k^2-6k+17\), \(m=8k-16\); consequently the finite bound is
\(\min(k^2-6k+17,16k-30)\). This is a finite rank statement, not an
estimate of a determinant or a large-degree return.
More explicitly, in block sizes \(1,1,m,m\), put
\(\mathcal G_d=[g_d,g_{d-1},F,T_d^*]\). Then the same formula is
\[
 HA_d-A_dH^*=\mathcal G_d
 \begin{pmatrix}
 0&\gamma_{d-1}^{-1}&0&0\\
 -\gamma_{d-1}^{-1}&0&0&0\\
 0&0&0&-M^{-1}\\
 0&0&M^{-1}&0
 \end{pmatrix}\mathcal G_d^*.
\]
Expanding its four nonzero blocks gives CD18 and CD20 exactly. The
matrix \(M\) depends only on the original coefficient frame, not on
the degree cutoff or on an inverse of the complete source covariance.

For a fixed cutoff the rank-two formula without the correction holds
if and only if \(L^*T_d^\perp\) is Hermitian, equivalently
\(FM^{-1}T_d=T_d^*M^{-1}F^*\). The stronger equation \(L=0\) is
sufficient at every cutoff and is equivalent to
\(Y^*\operatorname{im}V\subseteq\operatorname{im}V\).
For the original diagonal matrix \(Y\), whose eigenvalues are distinct,
this is equivalent to invariance under \(Y\): polynomial interpolation
at the finite distinct eigenvalues expresses \(Y^*\) as a polynomial
in \(Y\), and conversely expresses \(Y\) as a polynomial in \(Y^*\).
Neither invariance is asserted for the original OCF image merely from
its definition as an invariant-polynomial coefficient image. CD18--20
apply to its exact coefficients whether or not the displayed correction
vanishes.

All compression calculations also apply to CD17 with \(Y\) replaced
by its divided-jet multiplication matrix \(J\). For that possibly
nonnormal matrix, \(L=0\) means \(J^*\)-invariance. The diagonal
interpolation argument does not apply to a Jordan block, and no
equivalence with \(J\)-invariance is used.
In this general divided-jet statement \(r\) denotes the retained frame's
column count and \(m\) its complementary dimension. The OCF polynomial
rank formulas above are claimed only on the original simple-quartet
domain, not newly extended to a repeated-root coefficient construction.

\section{A finite Gamma example in which the correction is nonzero}

The following is an exact algebra fixture, not a sample of the original
period. It demonstrates why a general compression argument cannot
delete the term in CD18. Take the actual Gamma recurrence coefficient
\(a_1=1/2\), let \(M=M_\sigma\), and use
\[
 u=(0,i),\quad Y=\operatorname{diag}(0,i),\quad
 V=(1,i)^T,\quad Z=(1,-i)^T,\quad d=2.
\]
Then \(\gamma_0=M\), \(\gamma_1=M/2\),
\(f_0=(1,1)^T\), \(f_1=(0,i)^T\), and
\(f_2=(-1/2,-3/2)^T\). Direct multiplication gives
\[
 C_2=M^{-1}\begin{pmatrix}1&1\\1&3\end{pmatrix},\qquad
 A_2=4/M,\qquad W=2,\qquad H=i/2.
\]
Therefore \(HA_2-A_2H^*=4i/M\). But
\(g_1=1\), \(g_2=-1/2+3i/2\), so the compressed boundary term is
\[
 (g_2\overline{g_1}-g_1\overline{g_2})/\gamma_1=6i/M.
\]
The difference is the nonzero correction \(\mathcal E_2=2i/M\), with
the minus sign in CD18. This example has nonreal frame entries and
uses the retained positive Gamma mass. It establishes failure of an
uncorrected compression formula, not a claim about the value of the
original period's correction.

\section{Physical-coordinate displacement and mass powers}

Let \(X=cI+iY=\operatorname{diag}(v_\alpha)\),
\(h_j=i^jf_j\), and \(H_S=cI+iH\).
Multiplying CD16 by \(i\) gives the exact physical-coordinate identity
\begin{equation}\label{eq:S-displacement}
 XC_d+C_dX^*-2cC_d
 =\frac{h_dh_{d-1}^*+h_{d-1}h_d^*}{\gamma_{d-1}}.
 \tag{CD21}
\end{equation}
Indeed \(h_dh_{d-1}^*=if_df_{d-1}^*\) and
\(h_{d-1}h_d^*=-if_{d-1}f_d^*\). The corresponding original-frame
compression is
\begin{equation}\label{eq:S-compressed}
 H_SA_d+A_dH_S^*-2cA_d
 =\frac{(i^dg_d)(i^{d-1}g_{d-1})^*
       +(i^{d-1}g_{d-1})(i^dg_d)^*}{\gamma_{d-1}}
       -i\mathcal E_d.
 \tag{CD22}
\end{equation}
Every term is Hermitian, since \(\mathcal E_d\) is skew-Hermitian.
Thus both the anti-Hermitian displacement in the Gamma coordinate and
the shifted Hermitian displacement in the physical coordinate are
retained exactly.

As a diagnostic only, replacing \(d\sigma\) by \(a\,d\sigma\),
\(a>0\), leaves all monic polynomials, root coordinates, frames and
projection matrices unchanged. It multiplies every \(\gamma_j\) by
\(a\), so \(C_d,A_d,T_d,T_d^\perp,\mathcal E_d\) all scale by
\(a^{-1}\). Both sides of every CD formula scale identically.
The original source uses \(a=1\) and the exact mass in CD1--2.

\section{Verification and source boundary}

The companion checker compares finite polynomial sums against CD5--11,
the full and divided-jet displacement identities, the actual matrix
form of CD18--20 in declared nonreal algebra fixtures, the collision
pattern of the original quartet coordinate formula, the four cutoff
shifts, and the retained mass factors. Those fixtures do not evaluate
the actual period or change the original frame in the theorem.
The separate mathematical audit retains a second exact compression
derivation and its explicit nonzero-correction check.

\begin{thebibliography}{9}
\bibitem{AV} G. Akemann and G. Vernizzi,
\emph{Characteristic Polynomials of Complex Random Matrix Models},
arXiv:hep-th/0212051v2, original author source
\texttt{AkemannVernizzi.tex}, \texttt{CD} and its following footnote,
\url{https://arxiv.org/abs/hep-th/0212051v2}.
\bibitem{DLMF} T. H. Koornwinder, R. Wong, R. Koekoek and R. F. Swarttouw;
W. P. Reinhardt, \emph{Orthogonal Polynomials}, NIST DLMF Chapter 18,
original retained TeX for equations
\url{https://dlmf.nist.gov/18.19.E8},
\url{https://dlmf.nist.gov/18.19.E9a}, and
\url{https://dlmf.nist.gov/18.23.E7}.
The calculation above retains the exact density change \(dy=2dx\).
\bibitem{MR} Split-Zero programme,
\emph{The original observation kernel as an exact mixed-relation
determinant}, 20 September 2026, complete original MR1--23.

% BEGIN VERIFIED PUBLIC PROOF LINK
\href{https://github.com/KokunoYumeto/zeta-function-research-reader/blob/5992ad50c0f2a06921f1fcaded7b4e38097c0739/workbenches/splitzero-tandem/continuations/20260920-original-kernel-mixed-determinants/ORIGINAL_KERNEL_MIXED_DETERMINANT.tex\#L36}{Source proof MR1}.
% END VERIFIED PUBLIC PROOF LINK
\bibitem{OCF} Split-Zero programme,
\emph{An exact conductor frame and a fixed invariant-generator
recursion}, 14 September 2026, complete original OCF1--26,
especially OCF17 and OCF21--24.

% BEGIN VERIFIED PUBLIC PROOF LINK
\href{https://github.com/KokunoYumeto/zeta-function-research-reader/blob/5992ad50c0f2a06921f1fcaded7b4e38097c0739/workbenches/splitzero-tandem/continuations/20260920-original-kernel-mixed-determinants/providers/ORIGINAL_CONDUCTOR_STRIP_FRAME.tex\#L27}{Source proof OCF1}.
% END VERIFIED PUBLIC PROOF LINK
\bibitem{JM} Split-Zero programme,
\emph{The joint source in the original arithmetic metric: full
spectrum, observation, and heat return}, original JM21--28,
especially JM26--27.

% BEGIN VERIFIED PUBLIC PROOF LINK
\href{https://github.com/KokunoYumeto/zeta-function-research-reader/blob/5992ad50c0f2a06921f1fcaded7b4e38097c0739/workbenches/splitzero-tandem/continuations/20260920-original-kernel-mixed-determinants/providers/JOINT_MINIMUM_RECEIVERS.tex\#L286}{Source proof JM21}.
% END VERIFIED PUBLIC PROOF LINK
\bibitem{KG} Split-Zero programme,
\emph{Independent derivation of the complete Gamma reduction for the
original kernel}, 20 September 2026, KG1--21; and the accompanying
\emph{COERCIVITY\_INDEPENDENT}, C1--4, retaining the original DLMF
TeX sources and their norm conversion.

% BEGIN VERIFIED PUBLIC PROOF LINK
\href{https://github.com/KokunoYumeto/zeta-function-research-reader/blob/5992ad50c0f2a06921f1fcaded7b4e38097c0739/workbenches/splitzero-tandem/continuations/20260920-original-kernel-mixed-determinants/providers/KERNEL_GAMMA_INDEPENDENT.tex\#L73}{Source proof KG1}.
% END VERIFIED PUBLIC PROOF LINK
\end{thebibliography}
\end{document}
